% Mechanical relocation generated from the preservation ledger.
% Existing prose is included byte-for-byte from preserved blocks.

\section{Reduced Journey-Response Model}

Repeated evaluation of the complete event graph is unnecessary when routing
conditions are held fixed during demand estimation. The reduced model therefore
summarizes, for every retained OD-time cell, its feasible journeys, fixed
conditional shares, and expected contributions to observations. This section
defines that representation mathematically before the following subsection
illustrates its timetable search on a small network.

\subsection{Purpose and Assumptions}

The reduced model replaces repeated detailed assignment by a response for each
retained OD-time alternative.  It assumes that timetable feasibility and the
conditional distribution over scheduled journeys are computed before demand
estimation and remain fixed during that estimation.  Each response records the
expected contribution of one passenger journey to every modeled measurement.
Responses are additive, and frozen positive demand contributes a fixed offset.
Structural-zero cells have no response column and no estimated parameter.

This reduction is exact relative to its constructed journey choices when their
shares and the observation mapping remain fixed.  It is not a behavioral claim
that passengers use the same route throughout the day: every OD-time cell may
have a different feasible set and different shares.  Nor does it represent
endogenous crowding.  Approximation can enter through bounded transfer depth,
choice-set pruning, temporal aggregation, or a deliberately simplified demand
model; these assumptions must be distinguished from the algebra of the
response operator itself.

Journey construction and the initial shares use only the timetable, stop and
transfer topology, departure-time definition, generalized journey attributes,
and declared choice rules. Observed passenger counts are not used to select,
prune, or weight journeys. They enter only after the fixed response $B$ has
been constructed, through the likelihood used to estimate demand or gravity
parameters. This separation prevents the same counts from silently determining
both the response design and its fitted demand.

\subsection{Desired-Departure-Time Averaging}

A broad demand period must not be represented by whichever scheduled vehicle
happens to be convenient for preprocessing. Desired passenger departure time
is defined independently of supply. For OD-period cell $c=(o,d,t)$, the initial
model assumes
\[
 T_c^*\sim\operatorname{Uniform}(a_t,b_t).
\]
With $S_t$ equal-width midpoint samples,
\[
 \tau_{ts}=a_t+\left(s+\frac12\right)\frac{b_t-a_t}{S_t},
 \qquad w_{ts}=\frac1{S_t},
 \qquad s=0,\ldots,S_t-1.
\]
The timetable is queried only after $\tau_{ts}$ and $w_{ts}$ have been fixed.
Thus increased service frequency cannot itself attract more desired-time
weight. Uniform desired departures are an explicit assumption, not an
empirical conclusion.

Equal-count midpoints are retained for compatibility.  A fixed-step reference
instead partitions $[a_t,b_t]$ into intervals of prescribed maximum duration
and gives every midpoint its exact elapsed-duration weight, including a shorter
final interval.  Adaptive service-aware quadrature begins with coarse
subintervals, compares cached sparse journey or measurement responses, and
bisects only where feasibility, selected services, sparse support, or response
values change.  Timetable information selects numerical evaluation locations;
it never changes the uniform elapsed-time mass.  Sparse differences use a
normalized symmetric $ell_1$ measure over the union of supports.

Let $A_{c,s}$ be the measurement response generated from journeys feasible for
sample $s$. The averaged response is
\[
 \bar A_c=\sum_s w_{ts}A_{c,s}.
\]
For the default completed-journey interpretation, infeasible samples retain
their original mass in diagnostics, while the served response conditions on
the feasible set $F_c$:
\[
 p_c^{\mathrm{feasible}}=\sum_{s\in F_c}w_{ts},
 \qquad
 \bar A_c^{\mathrm{served}}
 =\frac{\sum_{s\in F_c}w_{ts}A_{c,s}}
        {p_c^{\mathrm{feasible}}}.
\]
This conditioning estimates completed public-transport journeys; it does not
estimate latent unserved demand and does not move an infeasible desired time to
an earlier, later, next-day, or supply-weighted departure.

For latent desired-departure demand, the alternative
``preserve mass'' policy uses the unconditioned response

\[
 \bar A_c^{\mathrm{latent}}=
 \sum_{s\in F_c}w_{ts}A_{c,s},
\]

so infeasible subintervals contribute an explicit zero response.  Conditional
route shares still sum to one, while a separate served-time fraction scales
the response atoms.  Feasible samples are never renormalized to replace the
missing mass.  The selected interpretation is part of artifact identity.

Adaptive diagnostics report evaluations, cache hits, accepted and refined
intervals, depth, feasible and infeasible time shares, support changes, a
heuristic relative response error, and unresolved probability mass.  The last
quantity is the weight stopped at a sample cap, depth limit, or minimum
resolution while still unstable; it is not advertised as a rigorous error
bound.  Quadrature changes preprocessing and cache identity but not the cost of
the subsequently compiled estimation objective.

The compatible pointwise implementation retains whole-period coverage and a
deterministic priority queue.  The recommended integral-response method instead
compares the probability-weighted parent midpoint integral with the sum of two
child-midpoint integrals.  Point supports need not agree.  Global refinement
selects the interval with the largest absolute contribution and stops when
\[
 \sum_I e_I\leq\epsilon_a+\epsilon_r\lVert Q\rVert_1.
\]
The integrated journey stage resolves active observations once and accumulates
the sparse expected measurement response, separated by destination.  Exact
trip-service identity remains a stricter diagnostic mode. Requested and
effective comparison modes are reported separately.

A bounded number of timetable departure and maximum-wait-window edges may seed
the initial numerical partition so that a narrow feasible-service window is
not silently missed.  Every interval retains elapsed-duration probability.
The timetable therefore informs quadrature locations but never supplies
behavioral desired-departure weights. Numerical integration error, unresolved
interval mass, and infeasible desired-time mass are distinct diagnostics.

Probability mass is canonicalized only at the aggregation boundary.  Finite
values within the configured tolerance of zero or one are projected to the
exact boundary and the raw value, canonical value, and correction are retained
in diagnostics.  Material excursions remain errors, and the strict
$0<p_c^{\mathrm{feasible}}\leq1$ invariant for retained choice sets is not
weakened.

The merger preserves the two probability layers
$P(s,r)=w_{ts}P(r\mid s)$. Identical scheduled paths or response atoms receive
the sum of their joint weights rather than new uniform shares. The desired-time
period identifies the OD cell, whereas actual boarding and alighting events
retain their timetable periods for observation matching. Consequently multiple
sample-specific searches still yield one response column per retained
OD-period cell.

The canonical model contract classifies each cell independently as free,
fixed-zero, or fixed-positive. Timetable feasibility never promotes a fixed
cell into the free parameterization. A feasible fixed-zero cell is retained
only in diagnostics and contributes neither a response column nor an offset.
A fixed-positive cell uses the same sampled averaging, remains outside the free
columns, and contributes its response multiplied by its fixed flow to the
measurement offset; absence of feasible sampled weight is a blocking error.
Only retained free cells define the origin-period groups requiring production
inputs. Searches use the exact declared sparse origin-period support rather
than an implicit Cartesian product of all origins and periods.

Zero feasible weight freezes a cell at zero. By default, positive feasible
weight below 0.5 excludes the cell from the initial model, weight from 0.5 to
below 0.9 retains it with a strong warning, and weight of at least 0.9 retains
it normally. These configurable thresholds are operational safeguards, not
universal behavioral constants. Sampling levels 1, 3, 6, 12, and higher allow
pre-estimation convergence analysis; one sample is a low-fidelity diagnostic,
not the recommended representation for long periods with vehicle-specific
counts.

\subsection{Timetable Feasibility and Journey Events}

For origin $o$, desired departure time $t$, and destination $d$, a bounded
round-based timetable search returns feasible journey alternatives
$j\in\mathcal J_{odt}$. A journey contains an ordered sequence of scheduled
vehicle legs. If it uses $L_j$ legs, it has $L_j-1$ transfers and contributes a
boarding and alighting event for every leg. Feasibility therefore depends on
the timetable, access window, minimum transfer time, directed footpaths, and
maximum transfer count.

Let $C_j$ be the generalized journey cost and let $\theta>0$ be the choice
dispersion. Conditional journey shares are
\[
 p_{j\mid odt}
 =\frac{\exp(-C_j/\theta)}
        {\sum_{k\in\mathcal J_{odt}}\exp(-C_k/\theta)}.
\]
Other declared pruning or dominance rules act before this normalization. A
different OD-time cell may have a different alternative set and different
shares even when its physical origin and destination are unchanged.

Here $t$ is the latent desired-departure period, not the period of the first
boarding. Every sampled query carries $t$ explicitly. Alternatives that first
board in different timetable periods are ranked, pruned, and normalized
jointly within the same $(o,d,t)$ choice set. Each boarding and alighting event
retains its realized period, so demand in period $t_0$ may affect measurement
rows in $t_1$. Period-semantics version~2 and sampling-integration schema~6
invalidate older sampled choices and downstream response artifacts.

For measurement $m$, let $h_{jm}$ be the number or weighted contribution of
events generated by journey $j$. The response of one passenger in cell
$c=(o,d,t)$ is
\[
 B_{mc}=\sum_{j\in\mathcal J_c}p_{j\mid c}h_{jm}.
\]
Hence $B_{mc}$ is an expected event contribution, not an observed route or an
individual simulated trajectory. Columns with identical event responses may be
represented by one equivalence class without changing forward or adjoint
products.
